
Machine learning reproducibility failures often surface hours into training when environment-stage API defects violate undocumented behavioral assumptions. This work presents a proof-of-concept validation of executable API contracts---lightweight validators that encode structural invariants (tensor shapes, dtypes), metamorphic properties (softmax probability sums, dropout identity), and composition requirements (cross-library consistency). We implement and evaluate three contract types on controlled synthetic defects. Structural contracts (N=2 real scenarios) detected shape and dtype mismatches at import time with zero false positives. Metamorphic contracts (N=50 synthetic scenarios) achieved 100\% detection of softmax sum violations and dropout identity failures in 3.7ms average execution time. Cross-framework validation contracts (N=30 synthetic test cases) detected dtype corruption, shape inconsistency, and numerical drift with 100\% detection rate and 0\% false positive rate. These proof-of-concept results demonstrate the technical feasibility of import-time and environment-stage behavioral validation, supporting the hypothesis that API contracts can detect certain classes of reproducibility defects before training begins. Further evaluation on larger defect corpora and production repositories is required to assess practical deployability and version stability.
